% QMB_n09_algorithmen.tex — Kapitel 9: Quantenalgorithmen
% Inhaltliche Grundlage: Dok. 073, 147, 176, 316, 343; eigene Darstellung
\chapter{Quantenalgorithmen: Deutsch, Grover, Shor und die Grenzen der Periode}
\label{ch:algorithmen}

\section{Die vier Standardalgorithmen}

Vier Algorithmen bilden das Lehrstück der Quantenrechnung. Sie zeigen,
was Quanteninterferenz leisten kann und was sie nicht leisten kann.

\subsection*{Deutsch: ein Bit in einer Abfrage}

Der Deutsch-Algorithmus beantwortet die Frage, ob eine unbekannte
Funktion $f:\{0,1\}\to\{0,1\}$ konstant oder balanciert ist, mit einer
einzigen Abfrage --- klassisch braucht man zwei.

In der Standard-QM: Man präpariert $|0\rangle$, wendet Hadamard an,
fragt das Orakel $U_f$ ab, wendet Hadamard an und misst. Ergebnis $0$:
konstant; Ergebnis $1$: balanciert. Der Trick ist das Phasenkickback:
$U_f|x\rangle = (-1)^{f(x)}|x\rangle$. Das Hadamard danach macht die
globale Phasendifferenz sichtbar.

In der FFGFT: Dieselbe Sequenz auf dem zylindrischen Phasenraum. Das
Orakel dreht die Phase des Torsionsvektors; Hadamard vertauscht $z$ und
$r$ und dreht um $\pi/2$. Das Messresultat ist identisch. Es gibt keinen
algorithmischen Unterschied --- nur der konzeptuelle Rahmen ist ein
anderer: Die Phase ist keine abstrakte komplexe Zahl, sondern die
Rotationsstellung der Torsionskonfiguration.\status{K}

\subsection*{Grover: Suche in $\sqrt N$ Schritten}

Grover sucht in einer unsortierten Datenbank mit $N$ Einträgen in $O(\sqrt N)$
Schritten statt $O(N)$ klassisch. Der Mechanismus: Man präpariert das
Register in einer gleichmäßigen Superposition aller $N$ Einträge,
markiert den gesuchten Eintrag durch Phasenkickback
($U_\mathrm{Grover}:|x\rangle\to(-1)^{\delta_{x,x_0}}|x\rangle$) und
spiegelt dann an der gleichmäßigen Superposition (Inversionsspiegel um
den Mittelwert). Jeder Grover-Schritt erhöht die Amplitude des gesuchten
Eintrags auf Kosten aller anderen. Nach $\pi/4\cdot\sqrt N$ Schritten
ist die Wahrscheinlichkeit des richtigen Eintrags nahe 1.

FFGFT-Einordnung: Grover-Iterationen sind in der FFGFT eine Folge von
Torsionsrotationen, die die Phase des gesuchten Eintrags systematisch
ausrichten. Die anderen Einträge landen in instabilen Konfigurationen;
die Auslesewechselwirkung treibt das System deterministisch in den
ausgerichteten Pol.\status{K} Die Vorhersagen sind identisch mit der
Standard-QM.

\subsection*{Shor: Faktorisierung in polynomialer Zeit}

Shors Algorithmus faktorisiert $N$ durch Periodenfindung. Sei $a$ mit
$\gcd(a,N) = 1$; die Funktion $f(x) = a^x\bmod N$ ist periodisch mit
Periode $r$. Die Quantenfouriertransformation (QFT) macht die Periode in
den Messwahrscheinlichkeiten sichtbar; die Kettenbruchnäherung extrahiert
$r$ aus dem Messwert. Aus $r$ folgen die Faktoren über $\gcd(a^{r/2}\pm1, N)$.

Zentrale Aussage des Korpus: Der Algorithmus zerfällt in zwei Vorgänge,
die nichts miteinander zu tun haben:\status{K}

\begin{center}
\small
\begin{tabular}{p{2cm}p{3cm}p{3cm}}
\toprule
 & \textbf{Aufbereitung} & \textbf{Transformation} \\
\midrule
Aufgabe & Modulare Exponentiation $a^x\bmod N$ & Periodenextraktion \\
Natur & arithmetisches Problem & Wellenoperation (Interferenz) \\
Aufwand & teuer ($O(n^3)$ Gatter) & billig ($O(n^2)$ Gatter) \\
\bottomrule
\end{tabular}
\end{center}

Die Aufbereitung macht $\approx95\,\%$ der Gatterkosten aus; die QFT nur
$\approx5\,\%$.\status{K} Shor spezifizierte nur den klassischen Prozess
$(a,1)\to(a,a^x\bmod N)$ --- die Implementierung der modularen
Exponentiation in Superposition ist das eigentliche Ingenieusproblem und
hängt nicht an quantenmechanischen Eigenschaften.

\section{$\phi$-QFT und Äquivalenzsatz}

Dok.~147 entwickelt eine alternative Fouriertransformation, die Phasen
$2\pi/\varphi^k$ statt $2\pi/2^k$ verwendet, wobei $\varphi$ der
Goldene Schnitt ist:
\begin{equation}
  \phi\text{-QFT}:|x\rangle\mapsto\frac{1}{\sqrt{Q_\varphi}}\sum_{y=0}^{Q_\varphi-1}e^{2\pi ixy/Q_\varphi}|y\rangle,
\end{equation}
mit $Q_\varphi = \varphi^n$ statt $Q = 2^n$.

\textbf{Hauptsatz:} Für Shors Algorithmus zur Faktorisierung von
$N < 2^{20}$ mit Fehlerwahrscheinlichkeit $\delta < 10^{-6}$ gilt:\status{K}
\begin{equation}
  P_{\mathrm{success}}(\text{Standard-QFT})\le P_{\mathrm{success}}(\phi\text{-QFT})\le P_{\mathrm{success}}(\text{Standard-QFT})+\xi.
\end{equation}
Die $\phi$-QFT ist zur Standard-QFT äquivalent, mit einer maximalen
Abweichung von $\xi\approx1{,}3\times10^{-4}$ --- vernachlässigbar.
Der Äquivalenzsatz zeigt: $\xi$ hat in $\mathbb{Z}/N\mathbb{Z}$ keine
strukturelle Funktion. Die Periode $r$ ist durch $\mathrm{lcm}(p-1,q-1)$
bestimmt, eine individuelle arithmetische Eigenschaft von $N$, nicht ein
universeller geometrischer Parameter.

\section{Die strukturellen Grenzen der Periode}

Drei Grenzen setzen dem, was Quantenalgorithmen leisten können, präzise
Schranken.

\textbf{1. Weyl-Gesetz und Zähldichte.}\status{K} Die Eigenwertdichte
jedes kompakten Zustandsraums wächst nach einem Potenzgesetz
$N(\lambda)\sim\lambda^{d/2}$. Die Primzahlzählfunktion $\pi(T)$ und die
Riemannschen Nullstellen wachsen dagegen wie $T\ln T$. Das ist keine
Potenz. Ein endlicher Zustandsraum kann ein solches Spektrum nicht
vollständig tragen. Die Weyl-Obstruktion: Ein $n$-Qubit-Register ist ein
kompakter Zustandsraum; arithmetische Zähldichten wachsen schneller als
jede kompakte Geometrie erlaubt.

\textbf{2. Aufbereitung als Flaschenhals.}\status{K} Die QFT macht
Periodizitäten sichtbar, die in den präparierten Zuständen bereits
vorhanden sind. Der limitierende Schritt ist die Zustandspräparation
--- die Auswertung von $a^x\bmod N$ --- nicht die Transformation selbst.
Kein geometrischer Parameter $\xi$ kann die arithmetische Komplexität
der Aufbereitung reduzieren: Sie ist eine Funktion der Primstruktur von
$N$, nicht der Geometrie des Trägers.

\textbf{3. Kein Quantenvorteil für RSA (Stand 2026).}\status{X} Ein
Vorteil muss auf der Shor-Faktorisierungskurve sichtbar werden: Der
Schwellenwert der Fehlerkorrektur muss demonstriert und gehalten werden;
die Aufbereitung muss in Superposition ausgeführt werden; das Ergebnis
muss mit RSA-relevanten Schlüssellängen (2048 Bit) repliziert werden.
Nichts davon ist Stand 2026 realisiert. Google Willow unterschreitet
zwar den Fehlerkorrektur-Schwellenwert für logische Qubits --- aber in
einem Regime weit unterhalb der für Shor nötigen Skalierung.

\section{Riemann-Zeta und FFGFT-Träger}

Dok.~316 und Dok.~343 klären die Beziehung zwischen der Riemann-Zeta
und dem FFGFT-Träger.

\textbf{Die Riemann-Zeta als Faktor der Torus-Spektral-Zeta.}\status{K}
Die geschlossene Form der $\mathbb{Z}^4$-Spektral-Zeta lautet:
\begin{equation}
  Z_{\mathbb{Z}^4}(s) = 8\,(1-4^{1-s})\,\zeta(s)\,\zeta(s-1).
\end{equation}
Das ist eine Identität. Die Nullstellen der Torus-Spektral-Zeta haben
drei Quellen: die Riemannschen Nullstellen (aus $\zeta(s)$), ihre
Verschobenen (aus $\zeta(s-1)$) und die Pol-Null $s = 1$. Die
Riemann-Hypothese ist damit eine Aussage über ein geometrisches Objekt
des Rahmens --- aber eine, die die Geometrie nicht beantworten kann:
Der Torus kann die Nullstellen als Faktor tragen, aber nicht erzeugen.

\textbf{$\mathbb{Z}_3$-Orbifold-Projektion.}\status{B} Die
$Z_3$-Projektion von $T^4/\mathbb{Z}_3$ transformiert die volle
Riemann-Zeta in eine Dirichlet-$L$-Funktion:
\begin{equation}
  \zeta_{T^4/\mathbb{Z}_3}(s) = L(s,\chi_0) = (1-3^{-s})\,\zeta(s).
\end{equation}
Das Euler-Produkt dieser $L$-Funktion ist nach der Eintrittstiefe
$k = \mathrm{ord}_p(3)$ hierarchisch geordnet; die physikalisch
sichtbaren Primzahlen $\{2,3,5,7,11,13\}$ sind genau jene mit $k\le6$.

\textbf{Galois-Struktur und Neutrinos.}\status{B} Die vier
Frobenius-Konjugationsklassen primitiver Elemente in $\mathrm{GF}(27)^*$
werden durch die FFGFT-Sektorpaarung $k\leftrightarrow-k$ zu zwei Klassen
gepaart: $\{f_3,f_4\}$ (Dirac, geladene Leptonen) und $\{f_1,f_2\}$
(Majorana, $\nu_1,\nu_2$ aktiv; $\nu_3$ in Dirac-Orbit $O_4$). Vorhersage:
$\nu_3$ hat keine Majorana-Phase.

\textbf{Weyl-Grenze.}\status{K} Keine kompakte Geometrie kann die
Nullstellen der Riemann-Zeta als Spektrum tragen. Die Nullstellendichte
$N(T)\sim T\ln T$ verlangt eine Dilatationsstruktur --- die ausgerollte
Domäne. Die Geometrie zeigt die Adresse an, liefert aber nicht die Antwort.
Drei naheliegende geometrische Zugänge sind ausgeschlossen: Keine
kompakte Geometrie liefert das $T\ln T$-Wachstum; $\xi$ verschiebt die
kritische Linie nicht; die $\xi$-Leiter diskretisiert die Nullstellen nicht.

\section{Deterministische Implementierung am PC}

Die Zustandsbrücke $(z,r,\theta)\leftrightarrow(\alpha,\beta)$ ist
bijektiv (Kapitel~\ref{ch:hilbertraum}). Daraus folgt unmittelbar:
Alle Logikbausteine --- Gatter, Algorithmen, Zustandsoperationen ---
sind als exakte Matrizenoperationen auf einem normalen PC
implementierbar. Es braucht kein Quantensubstrat, keine
Wiederholungen, keine Shot-Mittelung. Jedes Ergebnis ist in einem
einzigen Durchlauf vollständig bestimmt.

Dies wurde implementiert und mit IBM-QC-Ergebnissen verglichen
(\texttt{pruef\_logikbausteine.py}):

\begin{description}[leftmargin=2em,style=sameline,itemsep=0.3em]
\item[\textbf{Gatter}] $H$, $X$, $Z$, CNOT als exakte unitäre
  Matrizen; $X_{\mathrm{frak}}$ mit Drehwinkel $\pi K_{\mathrm{frak}}$,
  Abweichung $1{,}333\,\%$, Infidelität $4{,}4\times10^{-4}$.
  Ein Matrixprodukt --- kein Sampling.

\item[\textbf{Deutsch}] Konstantes und balanciertes Orakel: je
  ein einziger Lauf, exaktes Ergebnis.

\item[\textbf{Grover}] $n = 3$ bis $n = 10$ Qubits deterministisch;
  $P(\text{Ziel}) > 0{,}94$ in einem Durchlauf ohne Iteration.
  Skaliert auf normalem PC bis $n \approx 25$ (Speicher
  $2^{25}\times16$\,Byte $\approx 512$\,MB).

\item[\textbf{Bell-Zustand}] $P(00) = P(11) = 0{,}5$ exakt in
  einem Lauf. IBM-Kingston: Fidelität $0{,}9876$ über
  $50\times2048$ Shots. Der QC brauchte 100.000 Messungen;
  der PC: einen einzigen Durchlauf.

\item[\textbf{Periodenfindung}] $N = 15$ bis $N = 323$
  deterministisch, alle Faktoren korrekt. Kein Sampling,
  kein QC.\status{K}
\end{description}

Die deterministischen Ergebnisse stimmen mit den IBM-QC-Resultaten
innerhalb der Messgenauigkeit überein --- nicht annähernd, sondern
exakt. Der Unterschied liegt ausschließlich in der Methode: Der QC
misst probabilistisch und braucht viele Wiederholungen; der PC
berechnet die Amplitude direkt.

\section{Klassischer PC versus Quantenrechner}

\textbf{Deterministischer Einzeldurchlauf.}\status{K}
Die FFGFT-Zustandsabbildung $(z,r,\theta)\leftrightarrow(\alpha,\beta)$
ist bijektiv (Kapitel~\ref{ch:hilbertraum}). Daraus folgt: Alle
algorithmischen Schritte --- Periodenfindung, Faktorextraktion,
Gatteroperationen --- sind auf dem PC in einem einzigen Durchlauf
vollständig bestimmt. Kein Wiederholen, keine Shot-Mittelung, keine
Statistik. Der IBM-Kingston-Vergleich brauchte auf der QC-Seite
50~Läufe à 2048~Shots, weil der QC probabilistisch misst; die
FFGFT-Simulation am PC: ein einziger Lauf, identisches Ergebnis.

\textbf{Aufbereitung: PC schneller als QC.}\status{K}
Die modulare Exponentiation $a^x\bmod N$ ist auf dem PC direktes
klassisches Rechnen ohne Overhead. Der Quantenschaltkreis für
dieselbe Aufgabe ist durch Reversibilität, Uncomputation und
Toffoli-Gatter nachweislich \emph{teurer} als die klassische
Implementierung (Dok.~176). Shor hat die Aufbereitung in
Superposition nicht spezifiziert; sie ist und bleibt ein
arithmetisches Problem in jeder Technologie gleich aufwändig.

\textbf{Klassische Verfahren übertreffen die Periodensuche.}\status{K}
Die FFGFT-Periodensuche sequenziell: $O(N)$, gemessener
Wachstumsexponent der mittleren Ordnung $0{,}95$. Pollard~$\rho$
braucht 3--7 Schritte für dieselbe Aufgabe, Probeteilung 58--180.
Die klassische Resonanzsuche scheitert zudem doppelt: $O(N)$
Schritte \emph{und} eine Frequenzauflösung von $1/N^2$, die bei
RSA-Größen unterhalb jeder Zahlendarstellung liegt
(\texttt{s1\_periodensuche.py}, Dok.~176).

\textbf{Verschränkung trägt den Rechenvorteil nicht.}\status{K}
Gottesman-Knill (1998): Jeder Clifford-Schaltkreis ($H$, $S$,
CNOT) ist klassisch in Polynomzeit simulierbar --- unabhängig von
der erzeugten Verschränkung. Ein GHZ-Zustand mit $n = 1000$
Qubits ist maximal verschränkt und braucht klassisch nur $2n^2$
Bits Speicher. Was den Quantenvorteil trägt, ist ungeklärt;
Kontextualität und Nicht-Stabilisiertheit sind Kandidaten, keiner
ist etabliert.

\textbf{Die einzige Ausnahme: QFT in Superposition.}\status{K}
Der einzige echte Quantenvorteil liegt ausschließlich darin,
dass die QFT die Periode über alle $x$ gleichzeitig in
Superposition ausliest --- $O((\log N)^3)$ statt $O(N)$. Das
kann ein klassischer PC nicht replizieren. Dies ist die
\emph{einzige} Stelle, an der der QC schneller wäre.

\textbf{Derzeit nicht erreichbar.}\status{K}
Für RSA-2048 wären $4{,}1\times10^6$ physikalische Qubits und
$8{,}6\times10^9$ Gatter nötig --- eine Ressourcengrenze,
prinzipiell überwindbar, aber Stand~2026 um Größenordnungen
nicht erreicht. Google~Willow unterschreitet den
Fehlerkorrektur-Schwellenwert für logische Qubits, liegt aber
weit unterhalb der für Shor nötigen Skalierung. Die QFT-Ausnahme
existiert theoretisch; sie ist praktisch irrelevant.

\section{Experimentelle Validierung}

IBM Kingston (28. Mai 2026), 50 unabhängige Läufe, 2048 Shots pro Lauf:
Bell-Zustands-Fidelität $0{,}9876\pm0{,}0028$. Die Wiederholbarkeitsvarianz
ist statistisch verträglich mit gewöhnlichem Quanten-Shot-Noise (Verhältnis
beobachtete/erwartete Varianz: $1{,}02$). Eine frühere Auswertung von 3
Läufen hatte fälschlicherweise $40\times$ Determinismus gegenüber der QM
berichtet; das war Stichproben-Artefakt.\status{K}

Die $\xi$-Korrektur an CHSH-Werten liegt bei $\Delta\sim10^{-5}$ pro
Messung, weit unter NISQ-Rauschen. Die häufig zitierten kumulativen
Werte für $N = 73$ Qubits (Dok.~022 und Dok.~147) und die elementare
Abweichung $\xi/(2\pi)$ sind nicht direkt vergleichbar und dürfen nicht
gleichgesetzt werden (Kapitel~\ref{ch:chsh_exp}).

\quelle{Dok.~073 (Algorithmen in T0, 2025), Dok.~147 §3 ($\phi$-QFT-Äquivalenz), §4 (Bell-Tests), §5 (Shor), §8 (IBM-Validierung, Mai 2026), Dok.~176 (Aufbereitung vs.\ Transformation), Dok.~316 (Riemann, Weyl), Dok.~343 (Zeta-Galois)}
